%%%%%%%%%%%%%%%%%%%%%%%%%%%%%%%%%%%%%%%%%%%%%%%%%%%%%%%%%%%%%%%%%%%%%%
% Lớp 8 | Chương II. Hằng đẳng thức đáng nhớ và ứng dụng | Lập phương của một tổng. Lập phương của một hiệu
% Dạng 8D2?7-3: Tính nhanh, rút gọn biểu thức
%%%%%%%%%%%%%%%%%%%%%%%%%%%%%%%%%%%%%%%%%%%%%%%%%%%%%%%%%%%%%%%%%%%%%%


% nguồn: dan_tu_AI (nhập 16/11/2026)
\begin{bt}%[8D2H7-3]%[ma:ab80669e]
	Rút gọn biểu thức $A = (x+2)^3 - (x-2)^3$.
	\loigiai{
		Ta có:
		\begin{align*} 
			A &= (x^3 + 6x^2 + 12x + 8) - (x^3 - 6x^2 + 12x - 8) \\
			&= x^3 + 6x^2 + 12x + 8 - x^3 + 6x^2 - 12x + 8 \\
			&= 12x^2 + 16.
		\end{align*}
	}
\end{bt}

% nguồn: dan_tu_AI (nhập 16/11/2026)
\begin{bt}%[8D2H7-3]%[ma:ce463252]
	Rút gọn biểu thức $B = (2x+y)^3 - 8x^3 - y^3$.
	\loigiai{
		Ta có:
		\begin{align*} 
			B &= (8x^3 + 12x^2y + 6xy^2 + y^3) - 8x^3 - y^3 \\
			&= 8x^3 + 12x^2y + 6xy^2 + y^3 - 8x^3 - y^3 \\
			&= 12x^2y + 6xy^2.
		\end{align*}
	}
\end{bt}

% nguồn: dan_tu_AI (nhập 16/11/2026)
\begin{bt}%[8D2H7-3]%[ma:8744438a]
	Tính nhanh giá trị của biểu thức $C = 28^3 + 3 \cdot 28^2 \cdot 2 + 3 \cdot 28 \cdot 2^2 + 2^3$.
	\loigiai{
		Ta thấy biểu thức $C$ có dạng khai triển của lập phương một tổng:
		\begin{align*} 
			C &= 28^3 + 3 \cdot 28^2 \cdot 2 + 3 \cdot 28 \cdot 2^2 + 2^3 \\
			&= (28+2)^3 \\
			&= 30^3 \\
			&= 27\,000.
		\end{align*}
	}
\end{bt}

% nguồn: dan_tu_AI (nhập 16/11/2026)
\begin{bt}%[8D2H7-3]%[ma:430fc03b]
	Tính nhanh giá trị của biểu thức $D = 103^3 - 9 \cdot 103^2 + 27 \cdot 103 - 27$.
	\loigiai{
		Ta biến đổi biểu thức $D$ về dạng lập phương của một hiệu:
		\begin{align*} 
			D &= 103^3 - 3 \cdot 103^2 \cdot 3 + 3 \cdot 103 \cdot 3^2 - 3^3 \\
			&= (103-3)^3 \\
			&= 100^3 \\
			&= 1\,000\,000.
		\end{align*}
	}
\end{bt}

% nguồn: dan_tu_AI (nhập 16/11/2026)
\begin{bt}%[8D2H7-3]%[ma:9d7281a9]
	Rút gọn biểu thức $E = (x-1)^3 + 3(x-1)^2 + 3(x-1) + 1$.
	\loigiai{
		Xem $(x-1)$ là số thứ nhất và $1$ là số thứ hai, ta áp dụng hằng đẳng thức lập phương của một tổng:
		\begin{align*} 
			E &= (x-1)^3 + 3 \cdot (x-1)^2 \cdot 1 + 3 \cdot (x-1) \cdot 1^2 + 1^3 \\
			&= [(x-1) + 1]^3 \\
			&= x^3.
		\end{align*}
	}
\end{bt}

% nguồn: dan_tu_AI (nhập 16/11/2026)
\begin{bt}%[8D2V7-3]%[ma:54363dc7]
	Rút gọn biểu thức $M = (2x+1)^3 - 8x(x-1)^2 - 28x^2$. Từ đó tính giá trị của biểu thức $M$ khi $x = 0{,}5$.
	\loigiai{
		Khai triển và thu gọn biểu thức $M$:
		\begin{align*} 
			M &= (2x+1)^3 - 8x(x-1)^2 - 28x^2 \\
			&= (8x^3 + 12x^2 + 6x + 1) - 8x(x^2 - 2x + 1) - 28x^2 \\
			&= 8x^3 + 12x^2 + 6x + 1 - (8x^3 - 16x^2 + 8x) - 28x^2 \\
			&= 8x^3 + 12x^2 + 6x + 1 - 8x^3 + 16x^2 - 8x - 28x^2 \\
			&= (8x^3 - 8x^3) + (12x^2 + 16x^2 - 28x^2) + (6x - 8x) + 1 \\
			&= -2x + 1.
		\end{align*}
		Thay $x = 0{,}5$ vào biểu thức $M$ đã rút gọn, ta được:
		\[ M = -2 \cdot 0{,}5 + 1 = -1 + 1 = 0. \]
		Vậy khi $x = 0{,}5$ thì $M = 0$.
	}
\end{bt}
